\documentclass[11pt]{article}

\usepackage{amsmath,amssymb,amsthm}
\usepackage[margin=1in]{geometry}
\usepackage{mathtools}
\usepackage{stmaryrd}   % \llbracket \rrbracket
\usepackage{booktabs}
\usepackage[hidelinks]{hyperref}

%--- theorem environments -------------------------------------------------
\theoremstyle{plain}
\newtheorem{theorem}{Theorem}
\newtheorem{proposition}{Proposition}
\newtheorem{lemma}{Lemma}
\newtheorem{conjecture}{Conjecture}
\theoremstyle{definition}
\newtheorem{question}{Question}
\theoremstyle{remark}
\newtheorem{remark}{Remark}

%--- macros ---------------------------------------------------------------
\newcommand{\C}{\mathcal{C}}
\newcommand{\Set}{\mathbf{Set}}
\newcommand{\Vect}{\mathbf{Vec}}
\newcommand{\Vecfd}{\mathbf{Vec}_{\mathrm{fd}}}
\newcommand{\Fam}{\mathrm{Fam}}
\newcommand{\Poly}{\mathbf{Poly}}
\newcommand{\Cont}{\mathbf{Cont}}
\newcommand{\Gl}{\mathrm{Gl}}
\newcommand{\Add}{\mathrm{Add}}
\newcommand{\ext}[1]{\llbracket #1 \rrbracket}
\newcommand{\ihom}[2]{[#1,#2]}
\renewcommand{\lhd}{\vartriangleleft}
\newcommand{\op}{\mathrm{op}}
\newcommand{\pizero}{\pi_0}
\newcommand{\Hom}{\mathrm{Hom}}
\newcommand{\cod}{\mathrm{cod}}
\newcommand{\one}{\mathbf{1}}
\DeclareMathOperator{\Nat}{Nat}
\DeclareMathOperator{\Ext}{Ext}
\newcommand{\All}{\mathsf{All}}
\newcommand{\Exists}{\mathsf{Exists}}

\setlength{\parskip}{2pt}

%==========================================================================
\begin{document}

\begin{center}
{\Large\bfseries Containers over a base --- state of the program}\\[6pt]
{\normalsize A status memo}\\[10pt]
\begin{tabular}{rl}
\textbf{Author:} & MacBeth\\
\textbf{Date:} & 2026-09-03\\
\textbf{For:} & Neil Ghani (PI)\\
\end{tabular}
\end{center}

\noindent\rule{\textwidth}{0.4pt}\par
\begin{sloppypar}
{\small\emph{Corresponds to work-in-progress commit \texttt{53db0f7}; detailed write-ups:
\texttt{containers-over-a-base.pdf} (the approaches survey, with the three closedness theorems) and
\texttt{admissibility-composition-product.pdf} (the $\pizero$ characterisation) in
\texttt{scotmacbeth/work-in-progress}.}}
\end{sloppypar}
\noindent\rule{\textwidth}{0.4pt}

\medskip
\noindent Neil --- a fast catch-up on where the ``containers over different categories'' program
stands. The theory has crystallised into a clean two-paper package plus one live frontier. Sections
1--2 are the settled map and ledger; \S3--\S4 are the admissibility spine and the novelty gate;
\textbf{\S5 is the open-questions section you asked for and is the one to read}; \S6 corrects one point
in the thread you were replying to. I keep proofs to a one-line mechanism throughout --- everything
below is written up in full in the two PDFs.

%==========================================================================
\section{The question, and the map}

A \emph{container} in a category $\C$ is a pair (shapes, positions): a set $S$ of shapes together with
a position object $P_s\in\C$ for each $s$, extending to the endofunctor
$\ext{S,P}X=\coprod_{s}\ihom{P_s}{X}$. Over $\Set$ this is Abbott--Altenkirch--Ghani \cite{AAG}. The
moment one asks for ``a container over a base $\C\neq\Set$'', the notion forks into
three-or-four inequivalent constructions:
\begin{itemize}
\item[(1)] \textbf{External families}, $\Fam(\C^\op)$: an \emph{external} $\Set$-indexed coproduct of
internal homs, $\ext{S,P}=\coprod_s\ihom{P_s}{-}$. This is the only one that reaches additive bases
such as $\Vect$ (Dorta--Jarvis--Niu's $\Sigma\Pi\C$ construction \cite{DJN}).
\item[(2)] \textbf{Indexed / dependent-polynomial} $\Sigma\Pi\Delta$ over a locally cartesian closed
(LCC) base --- Gambino--Kock \cite{GambinoKock}.
\item[(3)] \textbf{Fibrational / comprehension}: positions as a fibration over shapes; subsumes (1) and
(2) as the family and codomain fibrations.
\item[(4)] \textbf{Walker's locally-subcartesian-closed} bases \cite{Walker}: an affine base-change
$\nabla_f\dashv\boxtimes_f$ (subpullbacks) replacing the top adjoint when the base is not LCC.
\end{itemize}

\paragraph{The organising insight.} These are not four rival definitions but a \emph{totally ordered
tower of weakenings} of the $\Sigma\Pi\Delta$ machinery, refining the local-cartesian-closure axis:
\[
\underbrace{\text{LCCC}}_{\text{G--K \cite{GambinoKock}}}\ \supset\
\underbrace{\text{pullbacks}}_{\text{Weber \cite{WeberPoly}}}\ \supset\
\underbrace{\text{subpullbacks}}_{\text{Walker \cite{Walker}}}\ \supset\
\text{spans}.
\]
LCC asks pullback to have a right adjoint globally; Weber's distributivity pullbacks weaken this to a
per-morphism (exponentiable-leg) condition; Walker weakens it further to an affine adjoint; at the
bottom sit bare spans with no base-change adjoint at all.

\paragraph{Where $\Vect$ falls.} $\Fam(\Vect^\op)$ falls off the \emph{bottom} of the tower: $\Vect$
fails every rung, always for the same two reasons. It is \textbf{non-extensive} --- the comparison
$A\amalg B\to A\oplus B$ from the disjoint union of underlying sets to the biproduct misses every
genuine linear combination, so $\coprod^{\Set}\subsetneq\oplus$ (coproduct $\neq$ biproduct). And it
has \textbf{no internal $\Pi$} --- it is not LCC, exponentiability fails. So \emph{only} approach (1)
can even form its semantics over $\Vect$. This is the honest reason the linear world needs the external
construction: not a preference, a forced move. (Von Glehn \cite{vonGlehn}, TAC 33 (2018), identifies
approach (2) with the codomain fibration and (1) with the family fibration; the two diverge exactly on
non-extensive / non-LCC bases.) The two axes --- extensivity and local cartesian closure --- are the
genuine coordinates of the whole landscape, and $\Vect$, failing both, is the discriminating example.

%==========================================================================
\section{Proved results: a ledger}

All four statements below are proved in full in \texttt{containers-over-a-base.pdf}. I give each as a
theorem plus its one-line mechanism. Write $\ext{-}\colon\Fam(\C^\op)\to[\C,\C]$ for the extension.

\begin{theorem}[T1, flagship --- fullness is unit-connectedness]
$\ext{-}$ is \emph{fully faithful} $\iff$ the monoidal unit $I$ is \emph{connected}, i.e.\
$\C(I,-)\colon\C\to\Set$ preserves small coproducts.
\end{theorem}
\noindent\emph{Mechanism:} the action on homs is $\prod_s$ of the comparison $\gamma\colon
\coprod_t\C(I,X_t)\to\C(I,\coprod_t X_t)$, and connectedness is exactly ``$\gamma$ always
bijective.'' \emph{This is not extensivity.} $\Set\times\Set$ is lextensive yet $\ext{-}$ is not full:
its unit $(1,1)$ is disconnected, so for source $(\{s\},(1,1))$ and target $(\{t_1,t_2\},((1,1),(1,1)))$
there are $2$ container morphisms against $4$ natural transformations --- the two ``crossed'' maps are
unrealised. This corrects the folklore that read ``extensivity'' off the $\Set$ case (where base and
value category coincide, fusing two different theorems --- ours, and Diers' reconstruction \cite{Diers}
which uses extensivity of the \emph{codomain}).

\begin{theorem}[T2 --- Dirichlet $\otimes$ closedness]
$\otimes$ on $\Fam(\C^\op)$ is closed $\iff$ the functor
$\Phi(Z)=\prod_t\coprod_r\C(M_r,Z\otimes Q_t)$ is familially representable.
\end{theorem}
\noindent\emph{Mechanism:} a right adjoint at $(T,Q),(R,M)$ represents a functor whose restriction to
the one-shape generators is $\Phi$; representability transfers through the free coproduct completion.
Closed on cartesian-closed bases (recovering the $\Poly$ Dirichlet hom, Niu--Spivak Ex.~4.78
\cite{NiuSpivak}) and on $\Fam_{\mathrm{fin}}(\Vecfd^\op)$ \emph{only}; it \textbf{fails} over both
$\Fam(\Vecfd^\op)$ and $\Fam(\Vect^\op)$ by \emph{dual} mechanisms --- infinitely many shapes break
product-closure, infinite-dimensional positions break single-factor representability. The single
load-bearing condition is simultaneous \emph{dualizability-and-summability}, met over $\Vect$ exactly
on the finite/finite-dimensional corner.

\begin{theorem}[T4-left --- $\lhd$ left-closedness is collapse]
The substitution product $\lhd$ is \emph{left}-closed $\iff$ it \emph{collapses}: positions tiny
($=$ dualizable $=$ finite-dimensional over $\Vect$) forces $\lhd=\otimes$, symmetric.
\end{theorem}
\noindent\emph{Mechanism:} tininess turns $\ihom{Z}{\coprod_t B_t}$ from the branching
$\coprod_{\tau\colon Z\to T}\prod_d B_{\tau(d)}$ (which makes the $\Set$ closure non-polynomial, hence
absent) into a genuine coproduct, whence $\lhd=\otimes$. Holds on
$\Fam_{\mathrm{fin}}(\Vecfd^\op)$. \textbf{The crown inversion:} extensivity, the property whose
\emph{failure} costs T1 and T2, is here the \emph{obstruction} --- non-extensivity is the \emph{repair}.
Extensivity $\perp$ $\lhd$-left-closedness. This is the survey's most surprising structural fact and is
why the classification is genuinely two-dimensional rather than a single ``goodness'' scale.

\begin{theorem}[Approach (3) anchored --- the logic of containers]
$\Cont(\cod)=\Fam(\cod^\op)$ is a \emph{bifibration} and a \emph{co-hyperdoctrine}. The fibre over
$(S,\{P_s\})$ is $\prod_s(\Set/P_s)^\op$; the quantifiers along the shape-collapse are the $\All/\Exists$
predicate liftings ($\All=\Pi$, $\Exists=\Sigma$).
\end{theorem}
\noindent\emph{Mechanism / dualisation:} the container hyperdoctrine is the \emph{fibrewise opposite}
of the standard $\Set$-family hyperdoctrine --- container-$\exists$ is $\All$, container-$\forall$ is
$\Exists$, $\wedge\leftrightarrow\vee$, and each fibre is a \emph{co-topos}. Consequently joint
Beck--Chevalley / Frobenius is \emph{one-sided}: the right-adjoint quantifiers carry it, the
left-adjoint ones fail, obstructed by co-topos non-distributivity. (Fibrational folklore --- $\Fam$
preserves fibrations \cite{Hermida}, $\cod$ is a bifibration \cite{Jacobs}; the fibrewise-op
presentation is von Glehn \cite{vonGlehn}.)

%==========================================================================
\section{The admissibility program (the spine of the generality question)}

This is where your ``for which bases does it generalise?'' question actually lives. Say $\C$ is
\emph{$\lhd$-admissible} if $\Fam(\C^\op)$ carries the composition product at all: for all $p,q$ there
is $p\lhd q$ with $\ext{p\lhd q}\cong\ext p\circ\ext q$. Per-object this is
\emph{absorptivity}: every object $X$ has $X\mapsto\ihom{P}{\ext q X}$ a coproduct of representables.
The whole story is in \texttt{admissibility-composition-product.pdf}.

\paragraph{Trichotomy.} The space of bases splits into three regimes.
\begin{itemize}
\item \textbf{Extensive pole.} admissible $\Rightarrow$ unit connected $\Rightarrow$ $\ext{-}$ fully
faithful and $(-)\lhd q$ has a \emph{left adjoint for every $q$}. One comparison map $\gamma$ powers
both fullness and the adjoint --- ``the same lemma applied twice.''
\item \textbf{Collapse pole.} admissible with unit $\cong 0$ (every object copower-tiny, $\lhd=\otimes$);
$(-)\lhd q$ has a left adjoint iff $|T|=1$.
\item \textbf{Inadmissible.} $\Set\times\Set$, $\Set_*$ (pointed sets), $\Gl((-)^2)$.
\end{itemize}

\paragraph{Extensive-CCC pole --- completely characterised (the headline).}
\begin{theorem}[$\pizero$-characterisation]
On a nontrivial infinitary-lextensive cartesian-closed base, $\Fam(\C^\op)$ is $\lhd$-admissible
$\iff$ the connected-components functor $\pizero$ preserves finite products.
\end{theorem}
\noindent This single criterion (Lemma~S: the shape object $[P,T\cdot\one]$ must be a copower of
$\one$) \emph{unifies the two inadmissible extensive bases into one phenomenon}: $\Set\times\Set$ fails
because its unit is disconnected ($\pizero(X\times Y)=X_1Y_1\sqcup X_2Y_2\neq(X_1\sqcup X_2)\times
(Y_1\sqcup Y_2)$), and $\Gl((-)^2)$ fails despite a \emph{connected} unit because
$\pizero(K\times K)=2\neq 1$ for the single edge $K$ (Weichsel's theorem on the graph tensor product
\cite{Weichsel}; $\Gl$ is a Grothendieck topos \cite{MacLaneMoerdijk,CarboniJohnstone}).

\paragraph{Connectedness is necessary but NOT sufficient.} $\Gl((-)^2)$ is the witness: a topos, a
connected unit, still inadmissible. So ``connected unit'' is strictly weaker than admissibility, and
the gap is exactly the external-vs-internal distributive law (internal composition of Gambino--Kock
polynomials holds in \emph{every} topos; the external-shape law is what $\pizero$-multiplicativity
controls).

\paragraph{Gap~1 is inhabited.} The region (admissible $\wedge$ non-collapse $\wedge$ non-cartesian) is
non-empty: $\Set\times\Vecfd$ is $\lhd$-admissible on its natural locus, with an explicit $\lhd$ and a
natural isomorphism --- the rigid $\Set$ factor forces the shape set $\coprod_s T^{A_s}$, the flexible
$\Vecfd$ factor absorbs it by biproduct collapse. Its unit is \emph{disconnected}, so the
extensive-pole extraction (admissible $\Rightarrow$ connected) genuinely needs its hypotheses:
admissibility alone does not force connectedness. The extensive pole is a \emph{sufficient} condition
for one-bit rigidity, not a characterisation of it.

%==========================================================================
\section{Novelty gate}

Dorta--Jarvis--Niu (arXiv:2305.05655, \cite{DJN}) build a $\otimes$ over a general base that matches
mine exactly under the identification $\Sigma\Pi\C=\Fam((\Pi\C)^\op)$. But $\Pi\C$ is the
\emph{coproduct-free} product completion, so the infinitary-lextensive-CCC hypotheses of the
admissibility theorems \emph{never hold there} --- DJN is therefore neither prior art for, nor a
counterexample to, the admissibility results. They prove no fullness/faithfulness result and pose the
base-condition question explicitly as future work (\cite[\S6]{DJN}). MacBeth's delta is: fullness (T1),
closedness (T2/T4), and change-of-enrichment (T3). \emph{One caveat, flagged:} their composition
product is \emph{weighted} and coincides with mine only at $\C=1$ (their direction object multiplies in
an outer factor $u_{i,a}$ that mine omits) --- the general comparison of $\lhd_{\mathrm{DJN}}$ with
$\lhd$ is open.

%==========================================================================
\section{Open questions (the part you asked for)}

\subsection*{THE FRONTIER --- Conjecture 6.2, the absorptive dichotomy}

\begin{conjecture}[Absorptive dichotomy]
Every absorptive object is a tensor of a \emph{copower-tiny} (flexible) part and a \emph{distributive}
(rigid) part. Hence every ``nice'' admissible base $=$ (rigid-extensive) $\times$ (flexible-additive),
and there is \textbf{no irreducible Gap-1 inhabitant}.
\end{conjecture}

\noindent A scoping result reduces the \emph{whole} conjecture to a single question:
\begin{question}
Is full (infinite-dimensional) $\Vect$ $\lhd$-admissible?
\end{question}
The reduction: any additive base is non-extensive, so it has no rigid factor; if it were admissible it
would be an \emph{irreducible} Gap-1 inhabitant, i.e.\ a counterexample to 6.2. So 6.2 holds
$\iff$ full $\Vect$ is \emph{inadmissible}. This in turn reduces to a homological question in the
additive functor category $\Add(\Vect,\Vect)$:
\begin{question}
Is $F=\Hom(k^{(\mathbb{N})},\bigoplus_{\mathbb{N}}-)=\prod_{\mathbb{N}}(\bigoplus_{\mathbb{N}}-)$
\emph{projective} in $\Add(\Vect,\Vect)$ --- equivalently, a coproduct of representables?
\end{question}

\noindent\textbf{Status: OPEN.} What is settled this cycle (``computed''), and what I want you to know
before spending time on it:
\begin{itemize}
\item The naive ``$\infty$-copower obstruction'' --- that $F$ fails to preserve coproducts/products,
hence cannot be a coproduct of representables --- is \textbf{not rigorous}. Every candidate separator
turns out to be a \emph{non-obstruction}: a coproduct of representables with infinite-dimensional
summands \emph{also} fails to preserve infinite coproducts/products. Preservation is simply the wrong
diagnostic here.
\item A Yoneda finite-row-support rigidity lemma (\textbf{Lemma R}) is \emph{proved}, but it kills only
one candidate summand, not the dichotomy.
\item The Baer--Specker / slenderness engine that refutes the $\mathbf{Ab}$-analogue \textbf{does not
transport}: over a field every vector space is free, so there are no slender objects to run the argument
on. (Cited in prose only --- see flag in the report.)
\item A rigorous resolution needs $\Ext^1(F,-)\neq 0$ in $\Add(\Vect,\Vect)$: honest
functor-category homological algebra over a field. This is the single highest-value open problem in the
program.
\end{itemize}
\emph{Honest current bet:} direction A (full $\Vect$ inadmissible $\Rightarrow$ dichotomy holds), by
analogy only --- no proof.

\subsection*{Secondary open questions}
\begin{itemize}
\item \textbf{Wildcard: nuclear / topological vector spaces} (where tiny $\neq$ finite-dimensional) ---
the one corner the extensive-pole screen could not close.
\item \textbf{Arity gap.} The closed convolutional tensors on $\Cont$ are $\otimes$ and $\triangleright_S$
(bounded arity); the infinite-arity case is open.
\item \textbf{Front C (``prompts vary'').} Does the branching stratification
$\lambda\text{-inv}\subsetneq\text{non-branch}\subsetneq\text{cartesian}\subsetneq\Pi\text{-Mendler}$
lift index-wise to the indexed-container monoid structures of De Pascalis--Uustalu--Veltr\`i
(arXiv:2509.25879, \cite{DPUV})? This is the self-contained fallback --- fully within reach without the
homological machinery.
\end{itemize}

%==========================================================================
\section{One correction to the thread you were replying to}

You asked whether $\delta\overset{?}{=}\Phi$ --- i.e.\ whether T2 closedness \emph{is}
Weber-distributivity. \textbf{Refuted; they are distinct axes.} Weber's $\delta$ tests the \emph{left}
position ($P_s$ exponentiable/tiny) --- that is exactly my T4-left collapse condition. T2's $\Phi$
tests the \emph{right} position and the target ($Q_t$ dualizable-and-summable) and \emph{never sees}
$P_s$. The two are two-way separable over $\Vect$ (a finite-dimensional-left, infinite-dimensional-right
datum makes $\delta$ invertible while $\Phi$ is unrepresentable, and the reverse datum does the
opposite). So \textbf{T2 does not fold into the LCC-weakening tower} --- it belongs to Weber's
\emph{other} machinery (parametric right adjoints / familial functors, \cite{WeberFamilial}), which
sits parallel to the tower, not on it. T4-left is the rung on the tower; T2 is off it.

%==========================================================================
\section*{Where I would go next}

The homological attack on projectivity of $F$ in $\Add(\Vect,\Vect)$ --- establishing
$\Ext^1(F,-)\neq 0$ over a field --- is the single highest-value open problem: it collapses the whole
absorptive dichotomy (Conjecture 6.2) to a yes/no, and it is the last thing standing between the
admissibility program and a complete classification. It is also genuinely hard, and none of the naive
separators work. Front C (the branching stratification lifting to indexed containers of \cite{DPUV})
is the self-contained fallback: fully within reach with the tools already built, and a clean
publishable increment if the homological front stalls. My plan is to open both --- a PROVE session on
$\Ext^1$ and an expository pass on the DPUV monoid structures --- and let whichever moves first set the
pace.

%==========================================================================
\begin{thebibliography}{99}

\bibitem{AAG} M.~Abbott, T.~Altenkirch, N.~Ghani.
\emph{Containers: constructing strictly positive types.}
Theoretical Computer Science \textbf{342} (2005), 3--27 (FoSSaCS 2003).

\bibitem{CarboniJohnstone} A.~Carboni, P.~T.~Johnstone.
\emph{Connected limits, familial representability and Artin glueing.}
Mathematical Structures in Computer Science / TAC (1995).

\bibitem{Diers} Y.~Diers.
\emph{Cat\'egories localisables.}
Th\`ese de doctorat d'\'etat, Universit\'e Paris VI, 1977. (Familial representability.)

\bibitem{DJN} S.~Dorta, A.~Jarvis, N.~Niu.
\emph{Monoidal structures on generalized polynomial categories.}
arXiv:2305.05655.

\bibitem{DPUV} G.~De Pascalis, T.~Uustalu, N.~Veltr\`i.
\emph{Monoid structures on indexed containers.}
arXiv:2509.25879.

\bibitem{GambinoKock} N.~Gambino, J.~Kock.
\emph{Polynomial functors and polynomial monads.}
Math. Proc. Cambridge Philos. Soc. \textbf{154} (2013), 153--192. arXiv:0906.4931.

\bibitem{Hermida} C.~Hermida.
\emph{Fibrations, logical predicates and indeterminates.}
PhD thesis, University of Edinburgh, 1993.

\bibitem{Jacobs} B.~Jacobs.
\emph{Categorical Logic and Type Theory.}
Studies in Logic and the Foundations of Mathematics \textbf{141}, North-Holland, 1999.

\bibitem{MacLaneMoerdijk} S.~Mac Lane, I.~Moerdijk.
\emph{Sheaves in Geometry and Logic.} Springer, 1992.

\bibitem{NiuSpivak} N.~Niu, D.~I.~Spivak.
\emph{Polynomial Functors: A Mathematical Theory of Interaction.}
Book draft / arXiv:2312.00990, 2023--2024.

\bibitem{vonGlehn} T.~von Glehn.
\emph{Polynomials and models of type theory.}
PhD thesis, University of Cambridge, 2015; see also Theory Appl. Categ. \textbf{33} (2018).

\bibitem{Walker} C.~Walker.
\emph{Locally subcartesian closed categories.}
arXiv:2607.10242.

\bibitem{WeberPoly} M.~Weber.
\emph{Polynomials in categories with pullbacks.}
Theory Appl. Categ. \textbf{30} (2015), 533--598. arXiv:1106.1983.

\bibitem{WeberFamilial} M.~Weber.
\emph{Familial 2-functors and parametric right adjoints.}
Theory Appl. Categ. \textbf{18} (2007), 665--732.

\bibitem{Weichsel} P.~M.~Weichsel.
\emph{The Kronecker product of graphs.}
Proc. Amer. Math. Soc. \textbf{13} (1962), 47--52.

\end{thebibliography}

\end{document}
